\documentclass[11pt,a4paper]{article}

\usepackage[margin=1.45cm,top=1.35cm,bottom=1.55cm]{geometry}
\usepackage[T1]{fontenc}
\usepackage[utf8]{inputenc}
\usepackage{lmodern}
\usepackage{microtype}
\usepackage{tabularx}
\usepackage{array}
\usepackage{xcolor}
\usepackage{enumitem}
\usepackage{fancyhdr}
\usepackage{lastpage}
\usepackage{hyperref}
\usepackage{ragged2e}

\definecolor{navy}{HTML}{1F2B3B}
\definecolor{lightgray}{HTML}{F4F4F4}
\definecolor{rulegray}{HTML}{D2D5D9}
\definecolor{textgray}{HTML}{4E5661}

\hypersetup{
    colorlinks=true,
    urlcolor=blue,
    linkcolor=navy
}

\setlength{\parindent}{0pt}
\setlength{\parskip}{4pt}
\renewcommand{\arraystretch}{1.18}
\setlist[itemize]{leftmargin=1.5em,itemsep=2pt,topsep=2pt}
\setlist[enumerate]{leftmargin=1.8em,itemsep=3pt,topsep=3pt}

\pagestyle{fancy}
\fancyhf{}
\lfoot{\small\color{textgray}AI tools map for a production AI / GenAI engineer role}
\rfoot{\small\color{textgray}Page \thepage}
\renewcommand{\headrulewidth}{0pt}
\renewcommand{\footrulewidth}{0pt}

\newcolumntype{L}[1]{>{\RaggedRight\arraybackslash}p{#1}}
\newcolumntype{Y}{>{\RaggedRight\arraybackslash}X}

\newcommand{\sectiontitle}[1]{%
    \vspace{0.35em}
    {\Large\bfseries\color{navy}#1}\par
    \vspace{0.25em}
}

\newenvironment{tooltable}
{%
    \noindent
    \begin{tabularx}{\textwidth}{|L{0.27\textwidth}|Y|}
    \hline
    \rowcolor{navy}
    \color{white}\bfseries Software / Platform / Framework &
    \color{white}\bfseries Two-line practical explanation \\
    \hline
}
{%
    \end{tabularx}
    \vspace{0.35em}
}

\begin{document}
    
    \begin{center}
        {\fontsize{23}{27}\selectfont\bfseries\color{navy}
        Software, Platforms, and Frameworks for the\\[2pt]
        AI Engineer Role}
    \end{center}
    
    \vspace{0.45em}
    
    \begin{center}
        \large\color{textgray}
        Target role: Production AI Engineer / AI Researcher focused on ML, GenAI,
        LLM systems, RAG, agentic workflows, Azure, Databricks, APIs, pipelines,
        deployment, monitoring, and business integration.
    \end{center}
    
    \vspace{0.7em}
    
    \sectiontitle{How to use this document}
    
    This map translates the job description into concrete tools you can mention
    in interviews, study selectively, or use to plan a small portfolio project.
    The strongest match for this specific role is the combination of
    \textbf{Python + Azure + Databricks + RAG/agent frameworks + MLOps/monitoring}.
    
    \sectiontitle{1. Production-ready AI applications}
    
    \begin{tooltable}
        \textbf{FastAPI} &
        Use it to expose ML or LLM functionality as clean REST APIs in Python.
        It is especially useful when a prototype must become a reliable internal
        service with typed inputs, validation, and automatic API documentation. \\
        \hline
        \textbf{Docker} &
        Package the application, model code, dependencies, and runtime into
        reproducible containers. This makes deployment to Azure, Kubernetes, or
        Databricks jobs more predictable and easier to maintain. \\
        \hline
    \end{tooltable}
    
    \sectiontitle{2. Traditional machine learning and data science}
    
    \begin{tooltable}
        \textbf{scikit-learn} &
        A strong default for classical ML, feature preprocessing, model selection,
        and baseline experiments. It is ideal for tabular classification, regression,
        clustering, and quick comparison before moving to heavier deep learning
        solutions. \\
        \hline
        \textbf{PyTorch} &
        A flexible deep learning framework for neural networks, custom training loops,
        and research-to-production workflows. It is useful when the task needs deep
        models, embeddings, sequence models, or integration with modern LLM tooling. \\
        \hline
    \end{tooltable}
    
    \sectiontitle{3. GenAI and LLM-based systems}
    
    \begin{tooltable}
        \textbf{Azure OpenAI Service} &
        Use it to access enterprise-managed LLMs through Azure security, networking,
        and governance. It fits internal assistant systems where authentication, data
        protection, monitoring, and integration with Microsoft services matter. \\
        \hline
        \textbf{Hugging Face Transformers} &
        Use it for open-source transformer models, embeddings, fine-tuning, and local
        or private experimentation. It is helpful when you need model control, domain
        adaptation, or alternatives to closed cloud APIs. \\
        \hline
    \end{tooltable}
    
    \newpage
    
    \sectiontitle{4. RAG systems}
    
    \begin{tooltable}
        \textbf{Azure AI Search} &
        Use it as the retrieval layer for enterprise RAG with keyword, vector,
        semantic, hybrid, and multimodal search patterns. It is a strong match when
        documents must be indexed securely and queried at scale inside Azure. \\
        \hline
        \textbf{LlamaIndex} &
        Use it to build document ingestion, indexing, retrieval, and query pipelines
        for RAG applications. It is especially good for connecting many data sources
        and turning unstructured internal documents into searchable LLM context. \\
        \hline
    \end{tooltable}
    
    \sectiontitle{5. Agentic AI and workflow automation}
    
    \begin{tooltable}
        \textbf{LangGraph} &
        Use it to build stateful, multi-step agents with controllable flows, tool calls,
        memory, retries, and human review points. It is better than a simple chatbot
        when the AI must plan, call tools, validate outputs, and continue across steps. \\
        \hline
        \textbf{Microsoft Copilot Studio} &
        Use it to create enterprise agents that connect to Microsoft 365, business
        data, and workflow actions. It is useful when business units need practical
        automation without every use case becoming a full custom software project. \\
        \hline
    \end{tooltable}
    
    \sectiontitle{6. AI services, APIs, and integration into existing systems}
    
    \begin{tooltable}
        \textbf{FastAPI} &
        Use it as the interface layer between models, databases, frontends, and
        enterprise systems. It lets you separate AI logic from application logic and
        expose stable endpoints for other teams. \\
        \hline
        \textbf{Azure Functions} &
        Use it for serverless AI-related tasks such as document processing, scheduled
        inference, event-based automation, or lightweight workflow triggers. It is
        useful when a full always-on service is unnecessary. \\
        \hline
    \end{tooltable}
    
    \sectiontitle{7. Modern AI architecture and model orchestration}
    
    \begin{tooltable}
        \textbf{LangChain} &
        Use it to connect prompts, tools, retrievers, memory, models, and output parsers
        into LLM applications. It is practical for quickly composing RAG and agent
        prototypes before hardening the most stable parts. \\
        \hline
        \textbf{Semantic Kernel} &
        Use it when building LLM applications in the Microsoft ecosystem with planners,
        plugins, and structured orchestration. It fits Azure-heavy environments where
        AI functions need to be connected to enterprise workflows and services. \\
        \hline
    \end{tooltable}
    
    \newpage
    
    \sectiontitle{8. Azure implementation}
    
    \begin{tooltable}
        \textbf{Azure AI Foundry} &
        Use it as the central workspace for building, evaluating, deploying, and
        managing AI apps and agents on Azure. It is suitable for enterprise GenAI
        projects that need models, agents, evaluations, safety checks, and deployment
        controls. \\
        \hline
        \textbf{Azure Machine Learning} &
        Use it for training, experiment tracking, model registration, online endpoints,
        batch jobs, and MLOps. It is a good fit for moving traditional ML and deep
        learning models from notebooks into governed production systems. \\
        \hline
    \end{tooltable}
    
    \sectiontitle{9. Databricks end-to-end AI solutions}
    
    \begin{tooltable}
        \textbf{Databricks Mosaic AI} &
        Use it for building, evaluating, deploying, and governing GenAI and ML
        applications on the Databricks Lakehouse. It fits roles where data engineering,
        model development, vector search, serving, and governance must live in one
        platform. \\
        \hline
        \textbf{Databricks Model Serving} &
        Use it to deploy ML models, foundation models, and agents behind scalable
        serving endpoints. It helps convert experiments into real-time or batch
        inference services that internal applications can call. \\
        \hline
    \end{tooltable}
    
    \sectiontitle{10. Data pipelines and feature pipelines}
    
    \begin{tooltable}
        \textbf{Apache Spark on Databricks} &
        Use it for distributed data processing, feature generation, ETL, and
        large-scale transformations. It is important when training data is large,
        comes from many systems, or must be processed regularly. \\
        \hline
        \textbf{Delta Lake} &
        Use it as a reliable storage layer for versioned, auditable, and transactional
        data pipelines. It helps keep training, inference, and analytics datasets
        consistent across teams and production jobs. \\
        \hline
    \end{tooltable}
    
    \sectiontitle{11. Training, inference, model registry, and monitoring}
    
    \begin{tooltable}
        \textbf{MLflow} &
        Use it to track experiments, package models, register model versions, and
        manage model lifecycle. For GenAI and agents, it also helps with tracing,
        evaluation, observability, prompt/version tracking, and production monitoring. \\
        \hline
        \textbf{Evidently AI} &
        Use it to monitor data drift, prediction drift, model quality, and regression
        risks after deployment. It is useful when production behavior may change
        because incoming data changes over time. \\
        \hline
    \end{tooltable}
    
    \newpage
    
    \sectiontitle{12. Performance, scalability, and maintainability}
    
    \begin{tooltable}
        \textbf{Kubernetes / Azure Kubernetes Service} &
        Use it to deploy, scale, and manage containerized AI services in production.
        It is useful when multiple services, APIs, workers, and model endpoints must
        run reliably with autoscaling and rollout control. \\
        \hline
        \textbf{GitHub Actions} &
        Use it for CI/CD: testing, linting, building Docker images, and deploying
        services after code review. It supports maintainability because every change
        can be checked and released through a repeatable pipeline. \\
        \hline
    \end{tooltable}
    
    \sectiontitle{13. Modern AI productivity tools}
    
    \begin{tooltable}
        \textbf{GitHub Copilot} &
        Use it as an AI coding assistant for Python, tests, refactoring, documentation,
        and code exploration. It can speed up implementation, but production code
        still needs review, security checks, and human ownership. \\
        \hline
        \textbf{Microsoft 365 Copilot} &
        Use it to accelerate everyday knowledge work such as summarizing documents,
        drafting specifications, and preparing meeting notes. In business-facing AI
        work, it can help translate stakeholder needs into clearer requirements and
        prototypes. \\
        \hline
    \end{tooltable}
    
    \sectiontitle{14. Collaboration with business units and rapid prototyping}
    
    \begin{tooltable}
        \textbf{Streamlit} &
        Use it to build quick internal demos, data apps, model dashboards, and
        proof-of-concept AI interfaces in Python. It helps business users understand
        and test an AI use case before expensive production engineering begins. \\
        \hline
        \textbf{Power BI} &
        Use it to communicate model outputs, KPI changes, and operational impact to
        non-technical stakeholders. It is useful when AI results must be connected to
        reporting, decision-making, and business ownership. \\
        \hline
    \end{tooltable}
    
    \newpage
    
    \sectiontitle{Suggested interview framing}
    
    A strong answer for this role should not list tools randomly. A clear production
    architecture would be: data and features in Databricks, MLflow for experiment
    tracking and model registry, Azure AI Foundry/Azure OpenAI for GenAI components,
    Azure AI Search or Databricks vector search for RAG, FastAPI for internal APIs,
    Docker/Kubernetes for deployment, and MLflow or platform monitoring for quality,
    latency, drift, and cost.
    
    \sectiontitle{Small portfolio project idea}
    
    Build an internal-document assistant: ingest PDFs or web pages, create embeddings,
    index them in Azure AI Search or Databricks vector search, answer questions with a
    RAG pipeline, expose the answer service through FastAPI, log traces and evaluation
    results with MLflow, and containerize the service with Docker. This single project
    demonstrates most requirements in the job advertisement.
    
    \sectiontitle{Priorities if time is limited}
    
    \begin{enumerate}
        \item Azure AI Foundry / Azure OpenAI + Azure AI Search for enterprise GenAI and RAG.
        \item Databricks + Spark + Delta Lake + MLflow for data, ML lifecycle, and deployment.
        \item LangChain or LlamaIndex for RAG and LangGraph for agentic workflows.
        \item FastAPI + Docker + Kubernetes/Azure services for production integration.
        \item GitHub Copilot and Microsoft Copilot only as productivity tools, not as substitutes for engineering knowledge.
    \end{enumerate}
    
    \sectiontitle{Sources consulted}
    
    The recommendations are based on the job description and current official
    documentation or product pages for the tools below. Links are included for
    verification and further study.
    
    \begin{itemize}
        \item Microsoft Foundry documentation:
        \url{https://learn.microsoft.com/en-us/azure/foundry/}
        \item Azure Machine Learning documentation:
        \url{https://learn.microsoft.com/en-us/azure/machine-learning/}
        \item Azure AI Search documentation:
        \url{https://learn.microsoft.com/en-us/azure/search/}
        \item Databricks Machine Learning documentation:
        \url{https://docs.databricks.com/en/machine-learning/}
        \item Databricks Model Serving documentation:
        \url{https://docs.databricks.com/en/machine-learning/model-serving/}
        \item LangChain documentation:
        \url{https://docs.langchain.com/}
        \item LlamaIndex documentation:
        \url{https://developers.llamaindex.ai/python/framework/}
        \item FastAPI documentation:
        \url{https://fastapi.tiangolo.com/}
        \item Kubernetes documentation:
        \url{https://kubernetes.io/docs/home/}
        \item MLflow documentation:
        \url{https://mlflow.org/docs/latest/}
        \item GitHub Copilot documentation:
        \url{https://docs.github.com/copilot}
        \item Streamlit documentation:
        \url{https://docs.streamlit.io/}
    \end{itemize}

\end{document}
